% Physical pages 3 and 4 of the concise study: the movement of the formulary,
% stated as an argument and not as a summary of what follows. The developed
% comparison of the three readings begins on the fifth page.

\section{The Propers: Themes and Movement}
\zlabel{triptych:concise:themes:start}

Two questions asked in the Temple in the last days before the Passion stand at
the centre of this Mass. A doctor of the law asks which commandment is the great one and is
given two for one; Jesus then asks the Pharisees whose son the Christ is, and
the psalm in which David calls his own descendant Lord ends the questioning.
Around that Gospel the Missal sets an apostle's plea from prison for a unity
kept in the bond of peace, a blessing on the nation whose God is the Lord, a
single cry for a hearing, three chants in which justice is confessed, mercy
asked and vows paid, and three orations that ask to escape the devil's
contagion, to be stripped of past and future sins, and to have our vices
cured. The Mass moves from a confession of God's justice at the door to the
payment of vows before him at the end, and the two questions stand at its
centre.

The Introit sets that opening note in one whole verse of the long alphabetical
psalm on God's law and in the first half of another. Its antiphon joins verse 137,
which opens the strophe \latin{Sade} on the justice of God, to the first half of
verse 124, from the strophe \latin{Ain}, in which the servant asks for help; the
words that complete it, ``and teach me thy justifications'', are not sung. The
confession of justice and the plea for mercy therefore stand side by side with
nothing between them. Two verses before verse 137, in a strophe the antiphon
passes over, the psalm prays \latin{Faciem tuam illumina super servum tuum}
(v.~135), ``Make thy face to shine upon thy servant'', with the verb and
construction the Offertory will use of the sanctuary. The Collect then asks
one gift under two infinitives: that God's people may shun \latin{diabolica
contagia} and follow him alone \latin{pura mente}, with a pure mind. Its verb
\latin{sectari}, a frequentative of \latin{sequi}, means to follow after
closely, to pursue; the first petition turns the heart from what corrupts it,
and the second turns it to the one it is to seek. Its \latin{mente} meets the
Gospel's \latin{in tota mente tua}.

The Epistle opens the moral part of the Letter to the Ephesians. Verses 1--3
describe the manner of life the calling asks --- humility, mildness, patience,
mutual support in charity, and care for the unity of the Spirit in the bond of
peace --- and verses 4--6 give its ground in a sevenfold \latin{unus}: one
body, one Spirit, one hope, one Lord, one faith, one baptism, one God and
Father of all. The Missal's \latin{Fratres} replaces the Vulgate's ``therefore'',
and its closing \latin{Qui est benedictus in saecula saeculorum. Amen}, which
is not part of Ephesians~4, gives the lesson a doxology of its own; verse 7,
which the Mass does not read, turns to the diverse gifts given to each. The
lesson stands where the Missal's course of Pauline epistles
puts it, after Eph 3:13--21 on the Sixteenth Sunday and before Ephesians
resumes on the Nineteenth, Twentieth and Twenty-first. The Gradual that follows
reverses its own psalm's order, taking the respond from verse 12 and the
versicle from verse 6, so that the blessing of the chosen nation is heard
before the word by which the heavens were made firm. The Alleluia gives the
day a single petition from the fifth penitential psalm, ``Hear, O Lord, my
prayer: and let my cry come to thee.'' The psalm moves from the afflicted
man's lament to God's enduring reign and the rebuilding of Sion: ``Thou shalt
arise and have mercy on Sion'' (v.~14).

The Gospel is the last of the controversies that run from Mt 21:23 to 22:46,
and after it ``neither durst any man from that day forth ask him any more
questions''. It has two halves. In the first a doctor of the law, tempting
him, asks for the great commandment and is answered with Deut 6:5 and, unasked,
with Lev 19:18 ``like to this''; in the second Jesus asks whose son the Christ
is and meets the answer ``David's'' with Ps 109:1 (Heb.\ 110:1). The Missal's
opening,
\latin{Accesserunt ad Iesum pharisaei}, drops the Vulgate's link with the
silenced Sadducees and lets the pericope begin with the Pharisees' approach.
The Creed follows the silence.

The Offertory answers the Gospel's silence with speech. It is a compilation and
not a continuous quotation: its narrative frame, \latin{Oravi Deum meum ego
Daniel, dicens}, stands outside the verses it cites and is formed from the
chapter's own first-person voice, and its petitions condense Dan 9:17--19.
The prayer they come from confesses the sins of Israel and the justice of God,
``To thee, O Lord, justice: but to us confusion of face'' (9:7).
Where the Vulgate reads \latin{ostende faciem tuam}, the chant has
\latin{illumina faciem tuam}, of a face made to shine. Its petitions leave out
the reason the chapter gives for daring to ask: ``for it is not for our
justifications that we present our prayers before thy face, but for the
multitude of thy tender mercies'' (9:18). The Secret then asks the
divine majesty that the holy things we do may strip us, \latin{exuant}, of past
sins and of future ones, and the Missal directs immediately after it the
Preface of the Most Holy Trinity. At the Communion the people is summoned to
pay its vows to the God who takes away the spirit of princes, in a psalm that
sings of the day ``When God arose in judgment, to save all the meek of the
earth'' (Ps 75:10); and the
Postcommunion asks that by God's sanctifying gifts our vices be cured and
eternal remedies come to us. The same Missal borrows that last clause for the
Collect of St John of God, where \latin{igne caritatis tuae} stands before it
and the cure is sought through the fire of God's charity, at a Mass that reads
this Sunday's Gospel.

The chants have a movement of their own: the servant first
confesses God's right judgment and asks for mercy; a chosen people blesses the
Lord whose word and spirit made the heavens; the poor man's cry seeks a
hearing; Daniel prays for the sanctuary and the people who bear God's name;
and those gathered round the Lord bring gifts and pay their vows. The Gospel
stands inside that movement. Its commandment tells the praying people what
an undivided heart is for, while its question about David's Lord tells them
whom they are confessing. The chants keep their own movement beside Matthew,
and the Gospel is heard inside the prayer that receives it.

The orations bind the chants and the Gospel into the act of this Mass. Before the
lessons, the Collect asks that a people be kept from what corrupts it and
follow the one God with a pure mind. After the Gospel, the Secret asks that the
holy things being done strip away sins already committed and guard against
those still to come. After Communion, the last prayer names those gifts as
sanctifying remedies that cure vice and reach into eternity. Thus the
formulary does more than place doctrine beside petition. It asks God to make
the worshipper capable of the love the Gospel commands, to preserve the unity
the Epistle says has already been given, and to heal the people who have
confessed both God's justice and their need of mercy.

Two habits run through the whole of it. The first is the oneness of God. The
Epistle says \latin{unus} seven times in three verses; the Collect asks that he
alone be followed, \latin{te solum Deum}; in the Gospel the Lord of David's
psalm sits at the right hand of the God whom the great commandment names; and
the Preface, the confession this Mass makes before the Canon, confesses him
\latin{unus es Deus, unus es Dominus}. The
second is the plea to be heard. The Introit asks God to deal with his servant,
the Alleluia asks that a cry reach him, and the Offertory asks that a prayer be
heard for a sanctuary and a people; between them stand the Gospel's questions,
the last of which nobody answers. What the Mass asks for, it asks three times
before it pays its vows.

Three readings of the whole Mass follow from these texts, each developed at
length in the expansive study. The first hears the formulary from the
Gospel's first half and the Epistle's appeal, as the prayer of a people
learning to love God wholly and to keep the unity that such love makes. The
second hears it from the Gospel's second half, the Epistle's ``one Lord \dots\
one God and Father'' and the Gradual's versicle, as the confession of the one
God in his Word and his Spirit, whose Son is David's Lord. The third hears it
from the Introit, the Offertory, the Alleluia and the Communion, as the prayer
of a humbled people that confesses God's justice and receives his mercy.
They differ in emphasis more than in what they identify: the first two divide
the Gospel between them and give the Mass its doctrine, while the third rests
on the chants and gives it its tone; and at several points they answer the
same question differently.
\zlabel{triptych:concise:themes:end}
